% =====================================================================
%  权利要求书（重构版）
%
%  主轴：证据身份跨表示保持不变 + 可重放证据账本 + 执行前状态门控。
%  具体摘要算法、阈值和文件布局均放入从属权利要求或实施方式，避免把
% 保护范围锁死在当前代码的某一组常量上。
% =====================================================================
\documentclass[12pt,UTF8,fontset=none]{ctexart}
\usepackage[claims]{cnipa}

\begin{document}

\claim{一种大语言模型智能体的工具调用控制方法，所述方法由智能体运行时执行，所述智能体运行时在一个用户回合内迭代地接收大语言模型生成的工具调用、执行所述工具调用以得到工具结果、将所述工具结果写入会话上下文并由所述大语言模型依据所述会话上下文生成下一次工具调用；}

\claimpara{其特征在于，所述方法还包括：}

\claimpara{S1、对每一次已完成的工具调用生成调用指纹和证据指纹，并在证据账本中写入与该次工具调用对应的观察记录；所述调用指纹由工具名称与经规范化处理的调用入参生成，所述证据指纹由工具输出的规范化表示生成，所述观察记录至少包括所述调用指纹、所述证据指纹或无证据标识、规范化规则版本、工具调用标识以及工具结果完成顺序；}

\claimpara{S2、生成所述证据指纹时，对工具输出进行确定性的结构规范化处理，并剔除预设的易变字段；当所述工具输出被变换为引用式结果、预览结果或上下文压缩后的恢复表示时，若所述变换后的表示携带所述工具输出的内容摘要值，则以所述内容摘要值作为证据身份依据，而不以存储路径、预览文本、调用标识或生成时间作为证据身份依据，使同一份工具观察在不同表示之间保持相同的证据指纹；}

\claimpara{S3、按照所述证据账本维护本用户回合的调用指纹计数、证据指纹集合、同结果计数和无新证据连续计数；当一次已完成的工具调用产生的证据指纹未出现在所述证据指纹集合中时，将其加入所述证据指纹集合并将所述无新证据连续计数清零；当未产生证据指纹、或产生的证据指纹已经存在于所述证据指纹集合中时，将所述无新证据连续计数累加；所述同结果计数用于表示同一调用指纹连续产生同一证据指纹的次数；}

\claimpara{S4、在执行大语言模型生成的一次新的工具调用之前，生成该次工具调用的调用指纹，并依据所述证据账本的状态进行执行前门控；当无新证据连续计数达到第一门控条件、所述调用指纹的历史计数达到第二门控条件且所述调用指纹对应的同结果计数达到第三门控条件时，拒绝执行该次工具调用并生成不代表外部工具观察的合成拒绝结果；所述合成拒绝结果不参与所述证据账本的更新，且调用指纹或工具能力发生改变的下一次工具调用仍被允许执行；}

\claimpara{S5、当所述会话上下文被压缩、恢复或重新加载时，从持久化的所述观察记录，或者从会话上下文中按调用标识配对的工具调用与工具结果，按照所述观察记录中记载的规范化规则版本和与 S1 至 S3 相同的确定性规则，重放并重建所述调用指纹计数、所述证据指纹集合、所述同结果计数和所述无新证据连续计数；所述重放过程不重新调用所述大语言模型或外部工具，并保持重放前后的证据新颖性判定关系一致。}

\claim{根据权利要求 1 所述的方法，其特征在于，步骤 S2 中所述结构规范化处理包括：将可解析的工具输出递归解析为结构化对象、对对象键名排序、保持数组元素顺序、按统一格式序列化，并从所述结构化对象中剔除时间戳、耗时、请求标识、链路追踪标识、心跳、更新时刻、开始时刻或结束时刻字段中的至少一项；对不可解析的工具输出合并连续空白字符。}

\claim{根据权利要求 1 所述的方法，其特征在于，当所述工具调用被标记为失败时，步骤 S2 中所述结构规范化处理包括：从工具输出中提取失败描述文本，将所述失败描述文本中的连续数字串替换为统一占位符并合并空白字符，以所得文本与工具名称共同生成所述证据指纹，使仅数字部分不同而错误类型相同的失败结果生成相同的证据指纹。}

\claim{根据权利要求 1 所述的方法，其特征在于，所述观察记录还包括工具输出的原始内容摘要值、规范化内容摘要值、表示类型、来源路径或来源窗口标识以及工具结果是否为合成结果的标记；所述证据账本仅将工具执行完成后获得的外部观察写入证据指纹集合。}

\claim{根据权利要求 1 所述的方法，其特征在于，步骤 S4 中所述第一门控条件、第二门控条件和第三门控条件由可配置策略确定，所述可配置策略至少包括关闭模式、仅记录与提示模式以及执行前阻断模式；在所述仅记录与提示模式下生成所述证据账本状态和策略变更指示而不拒绝工具调用。}

\claim{根据权利要求 1 所述的方法，其特征在于，所述智能体运行时读取工具能力描述，所述工具能力描述至少指示工具是否属于等待、状态查询或允许相同入参重复调用的轮询类工具；对于轮询类工具，将步骤 S4 中与无新证据连续计数对应的门控条件放宽为非轮询类工具对应条件的整数倍，或者采用所述工具能力描述指定的独立门控条件。}

\claim{根据权利要求 1 所述的方法，其特征在于，步骤 S2 还包括：当工具输出的原始长度或呈现长度超过根据当前上下文预算确定的外置阈值时，将工具输出原文写入位于预设产物目录内的产物文件，并以包含所述内容摘要值、所述产物文件的存储路径和有限预览的引用式结果替换会话上下文中的工具输出；当已存在具有相同内容摘要值且通过所述产物目录归属校验的产物文件时，复用所述产物文件。}

\claim{根据权利要求 7 所述的方法，其特征在于，所述有限预览由工具输出的首部片段、省略标记与尾部片段拼接得到，所述引用式结果还携带指示按行范围读取或按模式检索所述产物文件的检索提示；所述有限预览和所述检索提示不参与所述证据指纹的生成。}

\claim{根据权利要求 1 所述的方法，其特征在于，步骤 S5 中确定重放范围时，以会话上下文中最后一条非智能体运行时内部生成的用户消息之后的位置作为本用户回合的起点，以会话上下文末尾作为终点；遍历工具调用块并以调用标识建立调用入参记录，再以工具结果块的调用标识取回对应的调用入参；对于未找到对应调用入参的工具结果块使用预设的空入参结构进行重放。}

\claim{根据权利要求 1 所述的方法，其特征在于，步骤 S5 还包括：将由证据指纹集合的元素个数、已完成工具调用总次数、重复次数、空结果次数、当前无新证据连续计数以及规范化规则版本构成的聚合事实写入压缩后的摘要和跨压缩保留的固定上下文；所述聚合事实不包含工具输出原文。}

\claim{根据权利要求 1 所述的方法，其特征在于，当大语言模型在一次响应中生成多个并行工具调用时，所述证据账本仅依据已经完成的工具结果更新，尚未完成的工具调用不参与证据新颖性判定，且步骤 S4 的门控结果不撤销同一响应中已经派发执行的工具调用。}

\claim{根据权利要求 1 所述的方法，其特征在于，对每一次已完成的工具调用输出一条元数据事件，所述元数据事件包括工具名称、证据新颖性判定结果、证据指纹的截断值、无新证据连续计数、证据指纹集合的元素个数和已完成工具调用总次数，且所述元数据事件不包含工具输出原文。}

\claim{一种大语言模型智能体的工具调用控制系统，包括模型接口单元和工具执行单元，所述模型接口单元用于接收大语言模型生成的工具调用并向所述大语言模型提交会话上下文，所述工具执行单元用于执行所述工具调用以得到工具结果；}

\claimpara{其特征在于，所述系统还包括指纹生成单元、证据账本单元、执行前门控单元、表示变换单元、观察记录存储单元和上下文恢复单元；}

\claimpara{所述指纹生成单元用于生成调用指纹和证据指纹，所述表示变换单元用于在工具输出被外置、预览或压缩表示时保留内容摘要值，所述证据账本单元用于根据已完成工具结果更新所述调用指纹计数、证据指纹集合、同结果计数和无新证据连续计数，所述执行前门控单元用于在所述证据账本满足预设门控条件时在工具执行前生成不计入证据账本的合成拒绝结果，所述观察记录存储单元用于保存包括规范化规则版本的观察记录，所述上下文恢复单元用于按照相同规则重放所述观察记录以重建证据账本状态。}

\claim{一种计算机可读存储介质，其上存储有计算机程序，其特征在于，所述计算机程序被处理器执行时实现权利要求 1 至 12 中任一项所述的方法。}

\end{document}
